\section*{Chapter \ref{ch:g11:drift-to-spreading} --- From Continental Drift to Seafloor Spreading}

\begin{solution}{exo:g11:drift-to-spreading:1}
\omterm{def:g11:drift-to-spreading:spreading}{Seafloor spreading} is the making of new ocean floor at the crest of a
mid-ocean ridge, where molten rock rises and solidifies, followed by its
motion away from the ridge on both sides. The full rate is the speed at which
the two plates move apart; the half rate is the speed at which the floor on
one side moves away from the axis. For symmetric spreading, full rate $= 2
\times$ half rate.
\end{solution}

\begin{solution}{exo:g11:drift-to-spreading:2}
Above its \omterm{def:g11:drift-to-spreading:remanence}{Curie temperature}, about \qty{580}{\celsius}, magnetite cannot keep a
magnetisation of its own; only when it cools below it do its grains become
magnetised along the field and keep that direction. Heated again to
\qty{700}{\celsius}, above the \omterm{def:g11:drift-to-spreading:remanence}{Curie temperature}, the rock loses its old
record; as it cools again it records the field of that new time.
\end{solution}

\begin{solution}{exo:g11:drift-to-spreading:3}
The sound goes down and back: $d = 1500 \times 5.2/2 = \qty{3900}{m}$.
\end{solution}

\begin{solution}{exo:g11:drift-to-spreading:4}
The present normal interval began \num{0.78} million years ago: it has lasted
\num{780000} years. The Jaramillo event lasted $1.07 - 0.99 = \num{0.08}$
million years, \num{80000} years. Counting the boundaries younger than five
million years (0.78, 0.99, 1.07, 1.77, 1.95, 2.14, 2.15, 2.581, 3.04, 3.11,
3.22, 3.33, 3.58, 4.18, 4.29, 4.48, 4.62, 4.80, 4.89, 4.98): twenty reversals.
\end{solution}

\begin{solution}{exo:g11:drift-to-spreading:5}
$65\,000/40\,000 \approx 1.6$: the ridges would go about one and a half
times around the Earth.
\end{solution}

\begin{solution}{exo:g11:drift-to-spreading:6}
$20/0.78 \approx \qty{25.6}{km}$ per million years, that is a half rate of
about \qty{2.6}{cm} a year; the full rate is about \qty{5.1}{cm} a year.
\end{solution}

\begin{solution}{exo:g11:drift-to-spreading:7}
West: $45/3.58 \approx 12.6$~km per million years, \qty{1.26}{cm} a year;
east: $44/3.58 \approx 12.3$~km per million years, \qty{1.23}{cm} a year.
Full rate: $1.26 + 1.23 \approx \qty{2.5}{cm}$ a year. The two half rates
differ by about \qty{2}{\percent}: the spreading is very nearly symmetric.
\end{solution}

\begin{solution}{exo:g11:drift-to-spreading:8}
\qty{1.9}{cm} a year is \qty{19}{km} per million years: $990/19 \approx 52$
million years, a little older than the 49 million years of the oldest sediment
(about \qty{6}{\percent}).
\end{solution}

\begin{solution}{exo:g11:drift-to-spreading:9}
At \qty{20}{km} per million years: \qty{15.6}{km}, \qty{19.8}{km},
\qty{21.4}{km}, \qty{35.4}{km} and \qty{39.0}{km}. From the axis: normal from 0
to \qty{15.6}{km}, reversed to \qty{19.8}{km}, normal to \qty{21.4}{km}
(Jaramillo), reversed to \qty{35.4}{km}, normal to \qty{39.0}{km} (Olduvai),
then reversed.
\end{solution}

\begin{solution}{exo:g11:drift-to-spreading:10}
$64 \times 10^6/4.6 \times 10^9 \approx \num{0.014}$, about
\qty{1.4}{\percent} of the Earth's age. It supports Hess's argument that the
ocean floor is young because it is always being made and destroyed: no
ocean floor has survived from most of the Earth's history.
\end{solution}

\begin{solution}{exo:g11:drift-to-spreading:11}
At any moment the floor forms at the axis with the polarity of that moment,
and splits: half moves away on each side at the same speed. Two strips formed
at the same time therefore lie at the same distance from the axis, one on each
side, with the same polarity: each side is a copy of the other, reflected in
the axis. The symmetry is broken if the floor moves away faster on one side
than on the other (asymmetric spreading), or if the axis jumps to a new
position.
\end{solution}

\begin{solution}{exo:g11:drift-to-spreading:12}
(1) \qty{5}{cm} is $5 \times 10^{-5}$~\unit{km}: each year a strip
$65\,000 \times 5 \times 10^{-5} = \qty{3.25}{km^2}$ is added.
(2) $2.9 \times 10^8/3.25 \approx 9 \times 10^7$ years, about 90 million
years. (3) This is of the same order as the age of the oldest floor
(190 million years) and shorter than it: for the ocean floor to stay so
young, as much floor must be destroyed, somewhere, as is made at the ridges.
\end{solution}

\begin{solution}{exo:g11:drift-to-spreading:13}
The event lasted \num{0.08} million years. At \qty{12.5}{km} per million years
its stripe is $0.08 \times 12.5 = \qty{1}{km}$ wide; at \qty{75}{km} per
million years, $0.08 \times 75 = \qty{6}{km}$. Seen from the sea surface,
\qty{3}{km} above, a stripe much narrower than its depth produces only a weak
field, blurred together with the fields of its wide neighbours of the opposite
polarity; a stripe twice as wide as its depth stands out.
\end{solution}

\begin{solution}{exo:g11:drift-to-spreading:14}
Site 16: $176/11 \approx \qty{16}{m}$ per million years; site 20: $72/67
\approx \qty{1.1}{m}$ per million years. The sediment piles up at very
different rates from place to place and from time to time, so its thickness
does not give the age of the floor. The age is given by the fossils of the
oldest sediment, lying directly on the basalt.
\end{solution}

\begin{solution}{exo:g11:drift-to-spreading:15}
(1) West: $58/2.581 \approx \qty{22.5}{km}$ per million years,
\qty{2.25}{cm} a year; east: $70/2.581 \approx \qty{27.1}{km}$ per million
years, \qty{2.71}{cm} a year; full rate $128/2.581 \approx \qty{49.6}{km}$ per
million years, about \qty{5.0}{cm} a year. (2) About \qty{225}{km} west and
\qty{271}{km} east. (3) The two boundaries were together at the axis when the
reversal happened; since then they have moved apart at the full rate, however
the motion is shared between the two sides. Their separation divided by
their age is the full rate.
\end{solution}

\begin{solution}{pb:g11:drift-to-spreading:1}
\textbf{1.} Site 16 is the closest (\qty{191}{km}, 11 million years), site 20
the farthest (\qty{1270}{km}, 67 million years): the farther from the ridge,
the older the floor.

\textbf{2.} The sediment starts to settle on the basalt as soon as the basalt
has formed, so the oldest layer is about as old as the basalt. It can only be
a little younger, never older.

\textbf{3.} In kilometres per million years: site 16: $191/11 \approx 17.4$;
site 15: $380/24 \approx 15.8$; site 18: $506/26 \approx 19.5$; site 17:
$643/33 \approx 19.5$; site 14: $727/40 \approx 18.2$; site 19: $990/49
\approx 20.2$; site 20: $1270/67 \approx 19.0$.

\textbf{4.} $129.5/7 \approx \num{18.5}$~km per million years, about
\qty{1.85}{cm} a year.

\textbf{5.} Yes: their ratios (19.5) lie among those of the western sites.
The floor has spread at about the same rate on both sides: the spreading is
symmetric.

\textbf{6.} $t = d/u$: the age is proportional to the distance.

\textbf{7.} $u = 208\,276/10\,952 \approx \num{19.0}$~km per million years,
that is \qty{1.90}{cm} a year.

\textbf{8.} About $2 \times 1.9 = \qty{3.8}{cm}$ a year.

\textbf{9.} Yes: \qty{1.9}{cm} a year is slightly less than \qty{2}{cm} a
year.

\textbf{10.} $1270/19.0 \approx 67$ million years, the same as its oldest
sediment.

\textbf{11.} $191/9 \approx \num{21.2}$~km per million years,
\qty{2.1}{cm} a year.

\textbf{12.} $1270/71 \approx \num{17.9}$~km per million years,
\qty{1.8}{cm} a year.

\textbf{13.} $180\,329/9854 \approx \num{18.3}$~km per million years,
\qty{1.83}{cm} a year.

\textbf{14.} $(1.90 - 1.83)/1.83 \approx \qty{4}{\percent}$: the two rates
agree within the uncertainties of the ages and distances.

\textbf{15.} One age comes from fossils in the sediment, the other from the
magnetism of the basalt and the time scale of reversals. An error in one
method would not appear in the other: their agreement is a strong test of
\omterm{def:g11:drift-to-spreading:spreading}{seafloor spreading}.

\textbf{16.} $13.8 + 13.8 = \ang{27.6}$.

\textbf{17.} $111.2 \times \cos 30^\circ \approx \qty{96.3}{km}$ per degree;
$27.6 \times 96.3 \approx \qty{2660}{km}$.

\textbf{18.} $2660/19.0 \approx 140$ million years.

\textbf{19.} $2660/130 \approx \num{20.5}$~km per million years, about
\qty{2.0}{cm} a year, a little more than \qty{1.9}{cm} a year. Since the last
67 million years ran at about \qty{1.9}{cm} a year, the earlier opening was a
little faster: the floor between site 20 and the edge, $2660 - 1270 =
\qty{1390}{km}$, formed in $130 - 67 = 63$ million years, about
\qty{2.2}{cm} a year.

\textbf{20.} \textbf{The oldest ocean floor beside Africa at \ang{30} South
is about 130 million years old} (140 million years by extrapolating today's
rate): the South Atlantic began to open there in the Cretaceous period.
\end{solution}
